% Week 11: Revision, Peer Review, and Acting on Feedback
% BA-Eindwerkstuk Seminar - Leiden University
% Dr. Steven Denney

\documentclass[aspectratio=169,11pt]{beamer}

% ============================================
% Packages
% ============================================
\usepackage{tikz}
\usepackage{graphicx}
\usepackage{hyperref}
\usepackage{fontspec}
\usepackage{booktabs}

% Load custom theme
\usepackage{beamerthemethesis}

% ============================================
% Metadata
% ============================================
\title{BA-Eindwerkstuk Seminar}
\subtitle{Week 11: Revision, Peer Review, and Acting on Feedback}
\author{Dr. Steven Denney}
\institute{Korean Studies \\ Leiden University}
\date{May 8, 2026}

% ============================================
% Document
% ============================================
\begin{document}

% --------------------------------------------
% Title Slide
% --------------------------------------------
\begin{frame}[plain,noframenumbering]
    \titlepage
\end{frame}

% ============================================
\section{Opening}
% ============================================

\begin{frame}{Today's Agenda}
    \begin{enumerate}
        \item Where we are: from Assignment \#3 to the Final Manuscript
        \item Revision principles: re-seeing the argument
        \item Peer review activity: read and discuss a peer's empirical section
        \item The path to June 1
    \end{enumerate}
    \vspace{0.5cm}

    You submitted a 5{,}000-word empirical draft this week. Supervisor feedback comes in about a week; the Final Manuscript is due three weeks after that.
\end{frame}

\begin{frame}{Where We Are}
    \begin{table}
        \footnotesize
        \begin{tabular}{lll}
            \toprule
            \textbf{Milestone} & \textbf{Date} & \textbf{Status} \\
            \midrule
            Boot camp (Weeks 1--4) & Feb. 06 -- Feb. 27 & \textcolor{SageGreen}{Complete} \\
            Assignment \#1: Revised Proposal & Mar. 13 & \textcolor{SageGreen}{Complete} \\
            Assignment \#2: Preliminary Draft & Apr. 03 & \textcolor{SageGreen}{Complete} \\
            Peer review (Week 7) & Apr. 10 & \textcolor{SageGreen}{Complete} \\
            Ethics \& Protocol (Week 8) & Apr. 17 & \textcolor{SageGreen}{Complete} \\
            Assignment \#3: Empirical Draft & May 06 & \textcolor{SageGreen}{Submitted} \\
            \textbf{You are here} & \textbf{May 08} & --- \\
            Written feedback on A\#3 & within $\sim$1 week & \textcolor{Terracotta}{Coming} \\
            Final Manuscript & Jun. 01 & \textcolor{Terracotta}{In 24 days} \\
            \bottomrule
        \end{tabular}
    \end{table}
    \vspace{0.3cm}
    \begin{center}
        \itshape The remaining 24 days are for integrating, revising, and communicating the research you have already done.
    \end{center}
\end{frame}

% ============================================
\section{Revision Principles}
% ============================================

\begin{frame}{Revision is Not Editing}
    Most students conflate these. They are different operations, and the order matters.
    \vspace{0.15cm}
    \begin{table}
        \small
        \begin{tabular}{lll}
            \toprule
             & \textbf{Revision (global)} & \textbf{Editing (local)} \\
            \midrule
            Scope & Argument, structure, integration & Sentence, paragraph, citation \\
            \addlinespace
            Question & \textit{Does the thesis hold together?} & \textit{Is this prose clear?} \\
            \addlinespace
            Output & Sections moved, rewritten, cut & Cleaner sentences \\
            \addlinespace
            When & First & Last \\
            \bottomrule
        \end{tabular}
    \end{table}
    \vspace{0.3cm}
    Polishing a paragraph you may end up cutting is wasted effort. A weak link from evidence to claim is an argument problem dressed as a transition problem. Diagnose at the right level.
\end{frame}

\begin{frame}{Reverse Outlining}
    \footnotesize
    A reverse outline shows you what your draft \textit{actually does}, rather than what you intended.
    \vspace{0.1cm}

    \textbf{The method:}
    \begin{enumerate}
        \item For each paragraph, write a single sentence describing what it \textit{does} (e.g., ``introduces the second case,'' ``extends a counter-argument,'' ``summarizes the source'').
        \item Read those sentences as a list, in order. The structural problems will jump out.
    \end{enumerate}
    \vspace{0.05cm}
    \textbf{What to look for:}
    \begin{itemize}
        \item Two paragraphs doing the same thing $\rightarrow$ redundancy
        \item A claim with no paragraph supporting it $\rightarrow$ gap
        \item Out-of-order moves $\rightarrow$ restructure
        \item Paragraphs you cannot describe in one sentence $\rightarrow$ cut or rewrite
    \end{itemize}
    \vspace{0.15cm}
    \textbf{Try this:} take your introduction or empirical chapter, reverse-outline it (twenty minutes), then ask whether this is the argument you think you are making.
\end{frame}

\begin{frame}{Revising Against the Four Assessment Criteria}
    \footnotesize
    Your thesis is assessed on four criteria from the Protocol. Use them as a revision checklist.
    \vspace{0.1cm}
    \begin{table}
        \scriptsize
        \begin{tabular}{p{0.27\textwidth}p{0.65\textwidth}}
            \toprule
            \textbf{Criterion} & \textbf{Revision question} \\
            \midrule
            Knowledge and insight & Does the introduction frame a real, motivated question? Does the literature review show command of the field? \\
            \addlinespace
            Application of knowledge & Are methods justified? Is the empirical analysis well-executed and tied to the framework? \\
            \addlinespace
            Reaching conclusions & Does the discussion go beyond restating findings? Are limitations honest? \\
            \addlinespace
            Communication & Is the structure clear? Are citations correct? Does the prose serve the argument? \\
            \bottomrule
        \end{tabular}
    \end{table}
    \vspace{0.2cm}
    Most BA theses lose marks on \textbf{integration} (the empirical chapter does not answer the introduction's question) and on \textbf{conclusions} (the discussion just summarizes). Both are revision problems.
\end{frame}

% ============================================
\section{Peer Review and Feedback}
% ============================================

\begin{frame}{Peer Review: How This Will Work}
    \footnotesize
    \textbf{Format:} pairs (trios if odd number). \quad \textbf{Total time:} $\sim$35 minutes.
    \vspace{0.15cm}

    \textbf{Open the shared folder} and read your partner's empirical section from their most recent assignment. Five to eight pages is enough.
    \vspace{0.1cm}
    \begin{table}
        \scriptsize
        \begin{tabular}{lll}
            \toprule
            \textbf{Phase} & \textbf{Time} & \textbf{What happens} \\
            \midrule
            Setup & 5 min & I introduce the prompts; you pair up and open the section \\
            Silent reading & 10 min & You read your partner's section, annotating against three prompts \\
            Discussion (round 1) & 5 min & Reader A talks; Writer A listens and takes notes \\
            Discussion (round 2) & 5 min & Switch \\
            Whole-class debrief & 10 min & Share useful feedback; discuss as a group \\
            \bottomrule
        \end{tabular}
    \end{table}
    \vspace{0.1cm}
    \begin{alertblock}{The rule}
        The writer \textbf{listens without defending}. Take notes. Do not explain, do not push back, do not clarify. This is the same discipline you will need for written feedback from your supervisor.
    \end{alertblock}
\end{frame}

\begin{frame}{Three Reading Prompts}
    \footnotesize
    Annotate your partner's section against these three prompts. Write on the page or in a separate note.
    \vspace{0.15cm}

    \textbf{1. Argument trace.} In one sentence: what is the central claim of this section? Mark the sentence(s) where the claim appears, and the sentence(s) where the evidence supports it.
    \vspace{0.1cm}

    \textbf{2. Integration check.} Scan the introduction. Does this section visibly connect to the research question and the analytical framework laid out there? Where does the connection feel weak, missing, or assumed?
    \vspace{0.1cm}

    \textbf{3. One thing to keep, one thing to change.}
    \begin{itemize}
        \item Name a specific paragraph that works. Why?
        \item Name one specific \textit{structural} change you would suggest, not a prose tweak (e.g., ``move this section before the literature review,'' ``cut the second case,'' ``add a bridge paragraph between findings and discussion'').
    \end{itemize}
    \vspace{0.2cm}
    \textit{Structural feedback is more useful than prose feedback at this stage. Save line-editing for later.}
\end{frame}

\begin{frame}{Peer Review: Steps}
    \Large
    \vspace{0.4cm}
    \begin{enumerate}
        \setlength{\itemsep}{0.5cm}
        \item Silent reading (10 min)
        \item Discussion round 1 (5 min): Reader A speaks
        \item Discussion round 2 (5 min): Switch
    \end{enumerate}
    \vspace{0.6cm}
    \begin{center}
        \large \itshape The writer listens without defending.
    \end{center}
\end{frame}

\begin{frame}{Whole-Class Debrief}
    \large
    \vspace{0.4cm}
    \begin{center}
        \textbf{10 minutes.} Share useful feedback and discuss as a group.
    \end{center}
    \vspace{0.6cm}
    \begin{itemize}
        \setlength{\itemsep}{0.35cm}
        \item One useful piece of feedback you received
        \item A structural change you are now considering
        \item A question that came up during the activity
    \end{itemize}
\end{frame}

% ============================================
\section{Closing}
% ============================================

\begin{frame}{The Path to June 1}
    \begin{table}
        \small
        \begin{tabular}{ll}
            \toprule
            \textbf{Date} & \textbf{What happens} \\
            \midrule
            Within $\sim$1 week & Written feedback on Assignment \#3 \\
            Days 1--2 after feedback & Read once. 24-hour pause. Read again with the working method. \\
            Week 1 & Address Tier 1 (structure / argument) \\
            Week 2 & Address Tier 2 (evidence / citations); finish integration \\
            Week 3 & Address Tier 3 (prose / style); proofread; format \\
            \textbf{Jun. 01} & \textbf{Final Manuscript due} \\
            \bottomrule
        \end{tabular}
    \end{table}
    \vspace{0.3cm}
    Reach out to your supervisor between now and June 1 with anything you cannot resolve from the feedback alone. Office hours and meetings are still on.
\end{frame}

\begin{frame}{Key Takeaways}
    \begin{enumerate}
        \setlength{\itemsep}{0.35cm}
        \item \textbf{Revision is not editing.} Re-see the argument first. Polish prose last.
        \item \textbf{Reverse-outline at least one section} before the feedback arrives. You will find something.
        \item \textbf{Confusion is data.} When a reader misreads, fix the writing.
    \end{enumerate}
    \vspace{0.5cm}
    \begin{center}
        \footnotesize
        \textcolor{SlateBlue}{\href{https://scdenney.github.io/baks_thesis-seminar/}{Course website}} ~|~
        \textcolor{SlateBlue}{\href{https://scdenney.github.io/baks_thesis-seminar/protocol/}{Thesis Protocol}} ~|~
        \textcolor{SlateBlue}{\href{https://scdenney.github.io/baks_thesis-seminar/style-guide/}{Style Guide}} ~|~
        \textcolor{SlateBlue}{\href{https://scdenney.github.io/baks_thesis-seminar/assignments/}{Assignment \#3 Rubric}}
    \end{center}
    \vspace{0.3cm}
    \begin{center}
        {\color{AccentBlue}\large\bfseries Good luck with the final stretch.} \quad
        {\color{TextMuted}\itshape See you on the other side of June 1.}
    \end{center}
\end{frame}

\end{document}
</content>
</invoke>